%------------------------------------------------------------------------------%

\usepackage{apresentacao}
\usepackage{tabto}

\graphicspath{{./figs/}}

%------------------------------------------------------------------------------%
\newcommand{\mytitle}{Usos da Séries de Taylor}
\title   {\mytitle}
\author  {\large Luis Alberto D'Afonseca}
\date    {}

%------------------------------------------------------------------------------%
\begin{document}

%------------------------------------------------------------------------------%
\begin{frame}
    \titlepage
\end{frame}

%------------------------------------------------------------------------------%
\section{Séries binomiais}

%------------------------------------------------------------------------------%
\begin{frame}{Séries Binomiais}
Como avaliar funções potência?
\begin{itemize}
  \item \( f(x) = x^r \) \\[3mm]
  \item \( f(x) = \sqrt[3]{x} \) \\[3mm]
  \item \( f(x) = x^{\sfrac{5}{7}} \)
\end{itemize}
\end{frame}

%------------------------------------------------------------------------------%
\begin{frame}{Séries Binomiais}
Construir a série de Taylor da função \emph{\(\;f(x)=(1+x)^m\;\)} onde $m$ uma é constante
\begin{align*}
  \action<+->{f(x)       & = (1+x)^m                                \\[2mm]}
  \action<+->{f'(x)      & = (1+x)^{m-1}\; m                        \\[2mm]}
  \action<+->{f''(x)     & = (1+x)^{m-2}\; m(m-1)                   \\[2mm]}
  \action<+->{f^{(3)}(x) & = (1+x)^{m-3}\; m(m-1)(m-2)              \\}
  \action<+->{           & \;\;\vdots                               \\
              f^{(k)}(x) & = (1+x)^{m-k}\; m(m-1)(m-2)\cdots(m-k+1)}
\end{align*}
\end{frame}

%------------------------------------------------------------------------------%
\begin{frame}{Séries Binomiais}
Avaliando as derivadas em \emph{\(\;x=0\;\)} temos \(\;(1+x)^p=1\;\)
\begin{align*}
  f(0)       & = 1                        \\[2mm]
  f'(0)      & = m                        \\[2mm]
  f''(0)     & = m(m-1)                   \\[2mm]
  f^{(3)}(0) & = m(m-1)(m-2)              \\
             & \;\;\vdots                 \\
  f^{(k)}(0) & = m(m-1)(m-2)\cdots(m-k+1)
\end{align*}
\end{frame}

%------------------------------------------------------------------------------%
\begin{frame}{Séries Binomiais}
Os coeficientes da série são
\begin{align*}
	c_0 & = 1                        \\[2mm]
	c_1 & = m                        \\[2mm]
	c_2 & = \frac{m(m-1)}{2}         \\[2mm]
	    & \;\;\vdots                 \\
	c_n & = \frac{m(m-1)(m-2)\cdots(m-k+1)}{k!}
\end{align*}
\end{frame}

%------------------------------------------------------------------------------%
\begin{frame}{Séries Binomiais}
Para $m$ inteiro positivo

$c_k$ é o número maneiras de escolher $k$ elementos de um total de $m$

\vspace{\baselineskip}
\[
  \frac{m(m-1)(m-2)\cdots(m-k+1)}{k!} \;=\; \frac{m!}{k!(m-k)!} \;=\; \binom{m}{k}
\]
\end{frame}

%------------------------------------------------------------------------------%
\begin{frame}{Séries Binomiais}
Por similaridade definimos para $m$ real
\[
  \binom{m}{1} = m
\]
\hfill
\[
  \binom{m}{2} = \frac{m(m-1)}{2}
\]
\hfill
\[
  \binom{m}{k} = \frac{m(m-1)(m-2)\cdots(m-k+1)}{k!} \qquad k \geq 3
\]
\hfill
\end{frame}

%------------------------------------------------------------------------------%
\begin{frame}{Séries Binomiais}
Para \(-1<x<1\) temos
\[ (1+x)^m = 1 + \sum_{k=1}^{\infty} \binom{m}{k} x^k \]
\end{frame}

%------------------------------------------------------------------------------%
\section{Avaliando Integrais não Elementares}

%------------------------------------------------------------------------------%
\begin{frame}{Avaliando Integrais Não Elementares}
Como calcular a integral
\[
  \int \sin\left(x^2\right) \, dx
\]
\end{frame}

%------------------------------------------------------------------------------%
\begin{frame}{Avaliando Integrais Não Elementares}
Sabemos que
\[
  \sin\alpha = \alpha -\frac{\alpha^3}{3!} +\frac{\alpha^5}{5!} -\frac{\alpha^7}{7!} +\cdots
\]
\pause
fazendo \(\alpha = x^2 \)
\[
  \sin x^2 = x^2 -\frac{x^6}{3!} +\frac{x^{10}}{5!} -\frac{x^{14}}{7!} +\cdots
\]
\end{frame}

%------------------------------------------------------------------------------%
\begin{frame}{Avaliando Integrais Não Elementares}
Integrando termo a termo
\begin{align*}
  \action<+->{\int \sin x^2 \,dx}
  \action<+->{& = \int x^2
                - \frac{x^6}{3!}
                + \frac{x^{10}}{5!}
                - \frac{x^{14}}{7!}
                + \cdots  \,dx \\[5mm]}
  \action<+->{& = \int x^2 \,dx
                - \int \frac{x^6}{3!} \,dx
                + \int \frac{x^{10}}{5!} \,dx
                - \int \frac{x^{14}}{7!} \,dx
                + \cdots \\[5mm]}
  \action<+->{& = C
                + \frac{x^3}{3}
                - \frac{x^7}{7\times 3!}
                + \frac{x^{11}}{11\times 5!}
                - \frac{x^{15}}{15\times 7!}
                + \cdots}
\end{align*}
\end{frame}

%------------------------------------------------------------------------------%
\section{Avaliando Formas Indeterminadas}

%------------------------------------------------------------------------------%
\begin{frame}{Avaliando Formas Indeterminadas}
Calcular o limite
\[
  \lim_{x\to 1} \frac{\ln x}{\;x-1\;} \qquad\text{ indeterminação do tipo } \frac{0}{0}
\]

\pause
\vspace{\baselineskip}
Por l'Hôpital
\[
  \lim_{x\to 1} \frac{\ln x}{\;x-1\;} = \lim_{x\to 1} \frac{\;\sfrac{1}{x}\;}{1} = \frac{1}{1} = 1
\]
\end{frame}

%------------------------------------------------------------------------------%
\begin{frame}{Avaliando Formas Indeterminadas}
Solução alternativa usando Séries de Taylor

\pause
\vspace{\baselineskip}
Sabemos que
\[
  \ln(\alpha+1) = \alpha - \frac{\alpha^2}{2}
                         + \frac{\alpha^3}{3}
                         - \frac{\alpha^4}{4}
                         + \cdots
  \hspace{46mm} \abs{\alpha} < 1
\]

\pause
\vspace{0.5\baselineskip}
Substituindo $\alpha=x-1$ temos
\[
  \ln(x) = (x-1) - \frac{(x-1)^2}{2}
                 + \frac{(x-1)^3}{3}
                 - \frac{(x-1)^4}{4}
                 + \cdots
  \hspace{13mm} \abs{x-1} < 1
\]
\end{frame}

%------------------------------------------------------------------------------%
\begin{frame}{Avaliando Formas Indeterminadas}
Substituindo na função original
\begin{align*}
  \action<+->{\frac{\ln x}{\,x-1\,}} &
  \action<+->{= \frac{1}{\,x-1\,}
                \left( (x-1) - \frac{(x-1)^2}{2}
                             + \frac{(x-1)^3}{3}
                             - \frac{(x-1)^4}{4}
                             + \cdots
                \right)  \\[3mm]}
  \action<+->{& = 1 - \frac{x-1}{2} + \frac{(x-1)^2}{3} - \frac{(x-1)^3}{4} + \cdots}
\end{align*}
\null
\[
  \action<+->{  \lim_{x\to 1} \frac{\ln x}{\,x-1\,}}
  \action<+->{= \lim_{x\to 1} \left( 1 - \frac{x-1}{2}
                                       + \frac{(x-1)^2}{3}
                                       - \frac{(x-1)^3}{4}
                                       + \cdots
                              \right)}
  \action<+->{= 1}
\]
\end{frame}

%------------------------------------------------------------------------------%
\section{Identidade de Euler}

%------------------------------------------------------------------------------%
\begin{frame}{Identidade de Euler}
  Com os devidos cuidados todos os resultados sobre sequências e séries\\
  valem para os \emph{números complexos}
  \vspace{\baselineskip}
  \[
    z = a + bi \qquad \text{onde} \qquad a, b \in \R \qquad {\color{blue}i=\sqrt{-1}}
  \]
  \pause
  \vspace{0.5\baselineskip}
  \[
    i^2 = -1 \qquad i^3 = -i \qquad i^4 = 1 \qquad i^5 = i \qquad \cdots
  \]
\end{frame}

%------------------------------------------------------------------------------%
\begin{frame}{Identidade de Euler}
Definições Formais para $z\in\mathbb{C}$
\pause
\[
  \sin(z) = z - \frac{z^3}{3!} + \frac{z^5}{5!} + \frac{z^7}{7!} - \cdots
\]
\pause
\[
  \cos(z) = 1 - \frac{z^2}{2!} + \frac{z^4}{4!} + \frac{z^8}{8!} - \cdots
\]
\pause
\[
  e^z = 1 + z + \frac{z^2}{2!} + \frac{z^3}{3!} + \frac{z^4}{4!} +  \cdots
\]
\end{frame}

%------------------------------------------------------------------------------%
\begin{frame}{Identidade de Euler}
Escolhendo $z=i\theta$
\begin{align*}
  \action<+->{e^{i\theta}}
  \action<+->{& = 1 + (i\theta)
                    + \frac{(i\theta)^2}{2!}
                    + \frac{(i\theta)^3}{3!}
                    + \frac{(i\theta)^4}{4!}
                    + \frac{(i\theta)^5}{5!}
                    + \frac{(i\theta)^6}{6!}
                    + \cdots \\[2mm]}
  \action<+->{& = 1 +  i\theta
                    + \frac{i^2\theta^2}{2!}
                    + \frac{i^3\theta^3}{3!}
                    + \frac{i^4\theta^4}{4!}
                    + \frac{i^5\theta^5}{5!}
                    + \frac{i^6\theta^6}{6!}
                    + \cdots \\[2mm]}
  \action<+->{& = 1 +  i\theta
                    - \frac{\theta^2}{2!}
                    - i\frac{\theta^3}{3!}
                    + \frac{\theta^4}{4!}
                    + i\frac{\theta^5}{5!}
                    - \frac{\theta^6}{6!}
                    + \cdots \\[2mm]}
  \action<+->{& = \left( 1 - \frac{\theta^2}{2!}
                           + \frac{\theta^4}{4!}
                           - \frac{\theta^6}{6!}
                           + \cdots
                  \right)
                + i \left( \theta - \frac{\theta^3}{3!}
                                  + \frac{\theta^5}{5!}
                                  - \cdots
                    \right) \\[2mm]}
  \action<+->{& = \cos\theta + i\sin\theta}
\end{align*}
\end{frame}

%------------------------------------------------------------------------------%
\begin{frame}{Identidade de Euler}
\[
  e^{i\theta} = \cos\theta + i\sin\theta
\]
\pause
``Todas as Constantes da Matemática''
\[
  e^{i\pi} \;=\; \cos\pi + i\sin\pi \;=\; -1 + i0 \;=\; -1
\]
\[
  e^{i\pi} + 1 = 0
\]
\end{frame}

%------------------------------------------------------------------------------%
\section{Cálculo da Série de Taylor para o Arco Tangente}

%------------------------------------------------------------------------------%
\begin{frame}{Arco Tangente}
  Calcular a Série de Taylor de
  \[
    \arctan(x)
  \]

  \pause
  \vspace{\baselineskip}
  Sabemos que
  \[
    \frac{d}{dx}\arctan(x) = \frac{1}{1+x^2}
  \]
\end{frame}

%------------------------------------------------------------------------------%
\begin{frame}{Arco Tangente}
  Calculando as derivadas
  \begin{align*}
    \action<+->{\frac{d}{dx}     & \arctan(x) = \hpm \frac{1}{1+x^2}               \\[3mm]}
    \action<+->{\frac{d^2}{dx^2} & \arctan(x) =    - \frac{2x}{(1+x^2)^2}          \\[3mm]}
    \action<+->{\frac{d^3}{dx^3} & \arctan(x) = \hpm \frac{6x^2-2}{(1+x^2)^3}      \\[3mm]}
    \action<+->{\frac{d^4}{dx^4} & \arctan(x) =    - \frac{24x(x^2-1)}{(1+x^2)^4}}
  \end{align*}
  \action<+->{\emph{Não é um bom método!}}
\end{frame}

%------------------------------------------------------------------------------%
\begin{frame}{Arco Tangente}
  Cálculos formais (sem rigor)
  \pause
  \[
    \frac{d}{dx} \arctan(x) = \frac{1}{1+x^2}
  \]
  \pause
  Soma de uma série geométrica com $a=1$ e $r=-x^2$
  \[
    \frac{d}{dx} \arctan(x) = \frac{1}{1+x^2}
                            = \frac{a}{1-r}
                            = 1 - x^2 + x^4 - x^6 + x^8 - \cdots
  \]
  \pause
  Integrando os dois lados
  \[
    \arctan(x) = x -\frac{x^3}{3} + \frac{x^5}{5} - \frac{x^7}{7} + \frac{x^9}{9} - \cdots
  \]
\end{frame}

%------------------------------------------------------------------------------%
\begin{frame}{Problema}
  Como sabemos se os limites convergem?
\end{frame}

%------------------------------------------------------------------------------%
\begin{frame}{Demonstrando a Série de Taylor do Arco Tangente}
  \[
    \frac{1}{1+t^2} = 1 - t^2 + t^4 - t^6 + \cdots
                    + (-1)^n t^{2n}
                    + {\color<2->{blue}(-1)^{n+1} t^{2(n+1)} + \cdots}
  \]

  \pause
  \vspace{0.5\baselineskip}
  Usando a fórmula da soma da Série Geométrica a partir de $n+1$
  \[
    \frac{1}{1+t^2} = 1 - t^2 + t^4 - t^6 + \cdots
                    + (-1)^n t^{2n}
                    + {\color{blue}\frac{(-1)^{n+1} t^{2(n+1)}}{1+t^2}}
  \]

  \pause
  \vspace{0.5\baselineskip}
  Agora temos um número finito de termos
\end{frame}

%------------------------------------------------------------------------------%
\begin{frame}{Integrando a Soma Finita}
  Escrevendo a derivada como uma soma finita
  \[
    \frac{d}{dx} \arctan x = 1 - x^2 + x^4 - x^6
                           + \cdots
                           + (-1)^n x^{2n}
                           + \frac{(-1)^{n+1} x^{2(n+1)}}{1+x^2}
  \]

  \pause
  \vspace{0.5\baselineskip}
  Integrando
  \[
    \arctan x = x - \frac{x^3}{3}
                  + \frac{x^5}{5}
                  - \frac{x^7}{7}
                  + \cdots
                  + \frac{(-1)^n x^{2n+1}}{2n+1}
                  + \int_0^x \frac{(-1)^{n+1} t^{2(n+1)}}{1+t^2} \,dt
  \]
\end{frame}

%------------------------------------------------------------------------------%
\begin{frame}{A Série de Taylor Converge se o Erro Tende a Zero}
  \[
    R_n(x) = \int_0^x \frac{(-1)^{n+1} t^{2(n+1)}}{1+t^2} \,dt
  \]
  \pause
  como \(\;1+t^2\geq 1\;\)
  \[
    \abs{R_n(x)} \leq \int_0^{\abs{x}} t^{2(n+1)} \,dt \pause
      \; = \; \frac{t^{2n+3}}{2n+3}\evalat{0}{\abs{x}} \pause
      \; = \; \frac{\abs{x}^{2n+3}}{2n+3}
  \]
  \pause
  Portanto
  \[
    \lim_{n\to\infty} R_n(x) = 0 \qquad\text{quando}\qquad \abs{x} \leq 1
  \]
\end{frame}

%------------------------------------------------------------------------------%
\begin{frame}{Série de Taylor do Arco Tangente}
  Para todo \(\;x\in[-1,\,´1\,]\;\) temos
  \[
    \arctan x = x - \frac{x^3}{3}
                  + \frac{x^5}{5}
                  - \frac{x^7}{7}
                  + \frac{x^9}{9}
                  - \cdots
  \]
\end{frame}

%------------------------------------------------------------------------------%
\begin{frame}{\mytitle}
  \tableofcontents
\end{frame}

%------------------------------------------------------------------------------%
\end{document}
%------------------------------------------------------------------------------%
